% ==============================================================================
% Universidade Federal do Amazonas - UFAM
% Faculdade de Engenharia Elétrica e da Computação - FEEC
% Disciplina: Arquitetura de Sistemas Digitais e Laboratório
% Manaus, Amazonas, Brasil
%
% Slides para Apresentação em Vídeo-Aula (Segunda Entrega - MIPS Monociclo)
% Foco Central: Validação Rigorosa, Bateria de Testbenches e Algoritmo Sort
% ==============================================================================

\documentclass[aspectratio=169, 11pt]{beamer}

% Idioma e Tipografia
\usepackage[brazil]{babel}
\usepackage{fontspec}
\usepackage{amsmath, amssymb, amsfonts}
\usepackage{booktabs}

% Desativa símbolos de navegação do Beamer
\setbeamertemplate{navigation symbols}{}

% Tema Visual e Paleta Institucional UFAM
\usetheme{Madrid}
\usecolortheme{whale}
\usefonttheme{professionalfonts}

\definecolor{ufamblue}{RGB}{0, 43, 73}       % Azul Marinho Oficial UFAM
\definecolor{ufamlight}{RGB}{0, 102, 178}    % Azul Real Claro
\definecolor{ufamgold}{RGB}{196, 154, 69}    % Dourado Acadêmico
\definecolor{passgreen}{RGB}{34, 139, 34}    % Verde Sucesso

\setbeamercolor{structure}{fg=ufamblue}
\setbeamercolor{palette primary}{bg=ufamblue, fg=white}
\setbeamercolor{palette secondary}{bg=ufamlight, fg=white}
\setbeamercolor{palette tertiary}{bg=ufamblue!85!black, fg=white}
\setbeamercolor{titlelike}{parent=palette primary}
\setbeamercolor{block title}{bg=ufamblue, fg=white}
\setbeamercolor{block body}{bg=ufamblue!6!white, fg=black}
\setbeamercolor{block title example}{bg=ufamlight!90!black, fg=white}
\setbeamercolor{block body example}{bg=ufamlight!6!white, fg=black}

% Metadados da Apresentação
\title[MIPS Monociclo -- Validação e Testes]{Processador MIPS Monociclo com Jumps e Algoritmo Sort}
\subtitle{Segunda Avaliação: Implementação RTL, Bateria de Testbenches e Síntese em FPGA}

\author[Dercy, Vinicius]{%
  \textbf{José Dercy Vieira da Silva Filho} \and \\
  \textbf{Vinicius Menezes}%
}

\institute[UFAM / FEEC]{%
  Faculdade de Engenharia Elétrica e da Computação -- FEEC \\
  Universidade Federal do Amazonas -- UFAM
}

\date[Outubro 2026]{Arquitetura de Sistemas Digitais e Laboratório}

% ==============================================================================
% Início dos Slides
% ==============================================================================
\begin{document}

% ------------------------------------------------------------------------------
% Slide 1: Capa
% ------------------------------------------------------------------------------
\begin{frame}
  \titlepage
\end{frame}

% ------------------------------------------------------------------------------
% Slide 2: Objetivos da Segunda Entrega
% ------------------------------------------------------------------------------
\begin{frame}{Objetivos da Segunda Entrega}
  \begin{block}{Escopo do Projeto}
    Implementação e validação completa do processador \textbf{MIPS Monociclo de 32 bits} em Verilog HDL, sintetizável para a FPGA Intel MAX 10 (DE10-Lite), com foco na \textbf{verificação funcional exaustiva}.
  \end{block}

  \vspace{0.3cm}
  \textbf{Pilares Essenciais desta Entrega:}
  \begin{itemize}\setlength{\itemsep}{7pt}
    \item \textbf{Caminho de Dados com Jumps:} Suporte a desvios condicionais (\texttt{beq}, \texttt{bne}) e incondicionais (\texttt{j}, \texttt{jal}, \texttt{jr}).
    \item \textbf{Validação por Testbenches (Rubrica 3 -- Foco Principal):} Testes unitários para \textbf{todos} os módulos e teste de integração Top-Level executando o algoritmo real de ordenação (\texttt{sort.bin}).
    \item \textbf{Automação de Testes:} Bateria de regressão completa via Makefile e scripts bash (\texttt{make test}).
    \item \textbf{Síntese e Prototipagem em FPGA:} Análise, elaboração e bitstream gerados no Intel Quartus Prime.
  \end{itemize}
\end{frame}

% ------------------------------------------------------------------------------
% Slide 3: Visão Geral da Arquitetura
% ------------------------------------------------------------------------------
\begin{frame}{Visão Geral da Arquitetura MIPS Monociclo}
  O processador adota o \textbf{modelo Harvard} (barramentos separados para código e dados), estruturado em 4 subsistemas principais:

  \vspace{0.3cm}
  \begin{columns}[T]
    \begin{column}{0.48\textwidth}
      \begin{block}{1. Caminho de Dados (Datapath)}
        \small
        Registrador PC, somadores paralelos (PC+4 e Branch), extensores de sinal, multiplexadores em cascata, Banco de Registradores e ULA de 32 bits.
      \end{block}
      \vspace{0.2cm}
      \begin{block}{2. Unidade de Controle Integrada}
        \small
        Decodificação combinacional de \texttt{opcode} e \texttt{funct}, orquestrando todos os sinais de controle em um único ciclo de clock.
      \end{block}
    \end{column}

    \begin{column}{0.48\textwidth}
      \begin{exampleblock}{3. Memória de Instruções (ROM)}
        \small
        Armazenamento estático de 256 palavras de 32 bits inicializado via \texttt{\$readmemb} com o firmware compilado (\texttt{sort.bin}).
      \end{exampleblock}
      \vspace{0.2cm}
      \begin{exampleblock}{4. Memória de Dados (RAM)}
        \small
        Memória síncrona de 256 palavras para leitura (\texttt{lw}) e escrita (\texttt{sw}), onde o vetor desordenado é processado e ordenado *in-place*.
      \end{exampleblock}
    \end{column}
  \end{columns}
\end{frame}

% ------------------------------------------------------------------------------
% Slide 4: O Caminho de Dados e Lógica de Saltos
% ------------------------------------------------------------------------------
\begin{frame}{Caminho de Dados: Suporte Completo a Jumps}
  Para garantir chamadas e retornos de procedimentos (\texttt{jal} e \texttt{jr}) e desvios no mesmo ciclo, foi projetada uma \textbf{cascata de multiplexadores para o próximo PC}:

  \vspace{0.3cm}
  \begin{columns}[T]
    \begin{column}{0.50\textwidth}
      \textbf{Lógica de Decisão do PC:}
      \begin{enumerate}\setlength{\itemsep}{6pt}
        \item \textbf{Branch (\texttt{beq} / \texttt{bne}):}
          $$\text{pc\_src} = (\text{branch} \;\&\; \text{zero}) \;|\; (\text{bne} \;\&\; \sim\text{zero})$$
          Calcula: $(\text{PC}+4) + (\text{offset} \ll 2)$.
        \item \textbf{Jump Direto / JAL (\texttt{j} / \texttt{jal}):}
          $$\text{PC} \leftarrow \{\text{PC}+4[31:28],\; \text{target}[25:0],\; \text{2'b00}\}$$
        \item \textbf{Jump Register (\texttt{jr}):}
          $$\text{PC} \leftarrow \text{Reg}[rs] \quad (\text{normalmente } \text{Reg}[\$ra])$$
      \end{enumerate}
    \end{column}

    \begin{column}{0.48\textwidth}
      \begin{exampleblock}{Suporte a Funções e Link}
        \begin{itemize}\setlength{\itemsep}{6pt}
          \item \textbf{JAL (Chamada):} O endereço de retorno $\text{PC}+4$ é gravado no registrador \textbf{\$ra (\$31)} via multiplexador \texttt{mem\_to\_reg}.
          \item \textbf{JR (Retorno):} Desvia a execução de volta ao chamador.
          \item \textbf{Trava de Segurança:} O sinal \texttt{reg\_write} é mascarado como $\text{reg\_write} \;\&\; \sim\text{jr}$, evitando corrupção de registradores no retorno.
        \end{itemize}
      \end{exampleblock}
    \end{column}
  \end{columns}
\end{frame}

% ------------------------------------------------------------------------------
% Slide 5: Síntese e Prototipagem em FPGA (Rubrica 1)
% ------------------------------------------------------------------------------
\begin{frame}{Síntese e Prototipagem em FPGA (Rubrica 1)}
  \begin{block}{Configuração no Intel Quartus Prime}
    \begin{itemize}
      \item \textbf{Dispositivo Alvo:} FPGA \textbf{Intel MAX 10} (\texttt{10M50DAF484C7G}) na placa \textbf{Terasic DE10-Lite}.
      \item \textbf{Automação via Makefile:} Alvos \texttt{make quartus} (análise/elaboração) e \texttt{make quartus-build} (síntese, fitter e bitstream).
    \end{itemize}
  \end{block}

  \vspace{0.3cm}
  \begin{columns}[T]
    \begin{column}{0.48\textwidth}
      \begin{exampleblock}{Resultados de Compilação}
        \begin{itemize}\setlength{\itemsep}{4pt}
          \item \textbf{Análise e Elaboração:} \textbf{0 erros}.
          \item \textbf{Módulos RTL:} 15 módulos HDL reconhecidos e sintetizados com sucesso.
          \item \textbf{RTL Viewer:} Esquemático livre de latches espúrios e barramentos desconectados.
        \end{itemize}
      \end{exampleblock}
    \end{column}

    \begin{column}{0.48\textwidth}
      \begin{exampleblock}{Fluxo de Programação JTAG}
        \begin{itemize}\setlength{\itemsep}{4pt}
          \item \texttt{jtagconfig}: Reconhecimento do gravador USB-Blaster.
          \item \texttt{quartus\_pgm}: Gravação do arquivo \texttt{.sof} na SRAM do FPGA com comando único (\texttt{make quartus-pgm}).
        \end{itemize}
      \end{exampleblock}
    \end{column}
  \end{columns}
\end{frame}

% ------------------------------------------------------------------------------
% Slide 6: Estratégia de Testes e Bateria de Regressão (Rubrica 3)
% ------------------------------------------------------------------------------
\begin{frame}{Estratégia de Validação e Regressão (Rubrica 3 -- Foco Central)}
  \begin{block}{Filosofia de Verificação da Equipe}
    Nenhum módulo é integrado ao Top-Level sem antes ser exaustivamente validado por um testbench unitário dedicado, com verificação automatizada via assertivas (\texttt{if/else} com contadores de falhas).
  \end{block}

  \vspace{0.3cm}
  \textbf{Resumo da Bateria Completa de Regressão (\texttt{make test}):}
  \vspace{0.2cm}
  \centering
  \small
  \begin{tabular}{llcc}
    \toprule
    \textbf{Módulo Testado} & \textbf{Arquivo de Testbench} & \textbf{Qtd. Testes} & \textbf{Resultado} \\
    \midrule
    Unidade Lógica e Aritmética & \texttt{tb/tb\_alu.v} & 11 cenários & \textcolor{passgreen}{\textbf{100\% APROVADO}} \\
    Datapath Auxiliar (Adders, MUX, Extensor) & \texttt{tb/tb\_datapath\_aux.v} & 8 cenários & \textcolor{passgreen}{\textbf{100\% APROVADO}} \\
    Banco de Registradores (RegFile) & \texttt{tb/tb\_reg\_file.v} & 6 cenários & \textcolor{passgreen}{\textbf{100\% APROVADO}} \\
    Unidade de Controle Integrada & \texttt{tb/tb\_control\_unit.v} & 13 instruções & \textcolor{passgreen}{\textbf{100\% APROVADO}} \\
    Memória de Instruções (ROM) & \texttt{tb/tb\_instr\_mem.v} & 5 palavras & \textcolor{passgreen}{\textbf{100\% APROVADO}} \\
    Memória de Dados (RAM) & \texttt{tb/tb\_data\_mem.v} & 4 cenários & \textcolor{passgreen}{\textbf{100\% APROVADO}} \\
    Top-Level MIPS Monociclo (Sort Real) & \texttt{tb/tb\_mips\_monociclo\_top.v} & 5000+ ciclos & \textcolor{passgreen}{\textbf{100\% APROVADO}} \\
    \midrule
    \textbf{Total da Suíte de Testes} & \textbf{7 Testbenches Distintos} & \textbf{47+ Verificações} & \textcolor{passgreen}{\textbf{7 / 7 APROVADOS (0 FALHAS)}} \\
    \bottomrule
  \end{tabular}
\end{frame}

% ------------------------------------------------------------------------------
% Slide 7: Testes Unitários: Datapath Auxiliar e Banco de Registradores
% ------------------------------------------------------------------------------
\begin{frame}{Validação Unitária: Datapath Auxiliar e Registradores}
  \begin{columns}[T]
    \begin{column}{0.48\textwidth}
      \begin{block}{\texttt{tb\_datapath\_aux.v} (Componentes Auxiliares)}
        \small
        Valida os blocos isolados do caminho de dados:
        \begin{itemize}\setlength{\itemsep}{4pt}
          \item \textbf{PC Register:} Atualização síncrona com reset ativo em nível alto.
          \item \textbf{Somadores Paralelos:} Cálculo de $\text{PC}+4$ e endereço de branch com shift $(\text{offset} \ll 2)$.
          \item \textbf{Extensor de Sinal:} Imediatos positivos e negativos (ex: extensão de \texttt{0xFFFF} para \texttt{0xFFFFFFFF}).
          \item \textbf{Multiplexadores 2:1 e 3:1:} Seleção determinística de rotas de dados.
        \end{itemize}
        \textbf{Status:} \textcolor{passgreen}{\textbf{Todos os 8 testes aprovados com 0 erros.}}
      \end{block}
    \end{column}

    \begin{column}{0.48\textwidth}
      \begin{block}{\texttt{tb\_reg\_file.v} (Banco de Registradores)}
        \small
        Verifica as regras estritas da arquitetura MIPS:
        \begin{itemize}\setlength{\itemsep}{4pt}
          \item \textbf{Invariante \$zero (\$0):} Leitura fixa em 0 e bloqueio garantido contra qualquer escrita.
          \item \textbf{Leitura Assíncrona Dupla:} Leitura combinacional simultânea de duas portas (\texttt{rs} e \texttt{rt}).
          \item \textbf{Escrita Síncrona:} Escrita ativada por \texttt{reg\_write} sincronizada na borda de subida do clock.
          \item \textbf{Registrador \$ra (\$31):} Escrita e leitura validadas para suporte a \texttt{jal}.
        \end{itemize}
        \textbf{Status:} \textcolor{passgreen}{\textbf{Todos os 6 testes aprovados com 0 erros.}}
      \end{block}
    \end{column}
  \end{columns}
\end{frame}

% ------------------------------------------------------------------------------
% Slide 8: Testes Unitários: ULA e Unidade de Controle
% ------------------------------------------------------------------------------
\begin{frame}{Validação Unitária: ULA e Unidade de Controle}
  \begin{columns}[T]
    \begin{column}{0.48\textwidth}
      \begin{block}{\texttt{tb\_alu.v} (Unidade Lógica e Aritmética)}
        \small
        11 cenários de estresse numérico e lógico:
        \begin{itemize}\setlength{\itemsep}{4pt}
          \item \textbf{Lógicas:} \texttt{AND}, \texttt{OR}, \texttt{NOR}.
          \item \textbf{Aritmética com Sinal:} Soma e subtração com operandos positivos e negativos (ex: $10 + (-3) = 7$).
          \item \textbf{Zero Flag:} Ativação automática quando o resultado é nulo ($\sim|\,\text{alu\_result}$).
          \item \textbf{SLT com \texttt{\$signed}:} Validações com números com e sem sinal (ex: $-10 < 2 \rightarrow 1$).
          \item \textbf{Operação SLL:} Deslocamento à esquerda por \texttt{shamt} ($B \ll \text{shamt}$).
        \end{itemize}
        \textbf{Status:} \textcolor{passgreen}{\textbf{11 / 11 testes aprovados com 0 erros.}}
      \end{block}
    \end{column}

    \begin{column}{0.48\textwidth}
      \begin{block}{\texttt{tb\_control\_unit.v} (Unidade de Controle)}
        \small
        Decodificação exaustiva de todas as instruções:
        \begin{itemize}\setlength{\itemsep}{4pt}
          \item \textbf{Tipo-R:} \texttt{add}, \texttt{sub}, \texttt{and}, \texttt{or}, \texttt{slt}, \texttt{nor}, \texttt{sll}, \texttt{jr}.
          \item \textbf{Tipo-I:} \texttt{lw}, \texttt{sw}, \texttt{addi}, \texttt{slti}, \texttt{beq}, \texttt{bne}.
          \item \textbf{Tipo-J:} \texttt{j} e \texttt{jal}.
          \item \textbf{Trava Crítica do JR:} Confirmação de que \texttt{reg\_write = 0} durante a execução de \texttt{jr}, impedindo escrita indevida em registrador.
        \end{itemize}
        \textbf{Status:} \textcolor{passgreen}{\textbf{13 / 13 instruções validadas com 0 erros.}}
      \end{block}
    \end{column}
  \end{columns}
\end{frame}

% ------------------------------------------------------------------------------
% Slide 9: Testes Unitários: Subsistema de Memória (ROM e RAM)
% ------------------------------------------------------------------------------
\begin{frame}{Validação Unitária: Memória de Instruções e Dados}
  \begin{columns}[T]
    \begin{column}{0.48\textwidth}
      \begin{block}{\texttt{tb\_instr\_mem.v} (Memória de Instruções -- ROM)}
        \small
        Valida a leitura do código de máquina montado:
        \begin{itemize}\setlength{\itemsep}{5pt}
          \item \textbf{Alinhamento por Palavra:} Indexação por $\text{addr}[9:2]$, ignorando os 2 bits menos significativos (divisão por 4 em hardware).
          \item \textbf{Carga do Firmware (\texttt{sort.bin}):} Leitura exata das instruções binárias nos endereços \texttt{0x0}, \texttt{0x4}, \texttt{0x8}, \texttt{0xC} e \texttt{0x10}.
          \item \textbf{Leitura Combinacional:} Resposta imediata sem latência de ciclo de clock.
        \end{itemize}
        \textbf{Status:} \textcolor{passgreen}{\textbf{Leituras de palavras aprovadas (0 erros).}}
      \end{block}
    \end{column}

    \begin{column}{0.48\textwidth}
      \begin{block}{\texttt{tb\_data\_mem.v} (Memória de Dados -- RAM)}
        \small
        Valida a integridade do armazenamento de dados:
        \begin{itemize}\setlength{\itemsep}{5pt}
          \item \textbf{Escrita Síncrona (\texttt{sw}):} Dados gravados com segurança apenas na borda de subida do clock quando \texttt{mem\_write = 1}.
          \item \textbf{Proteção de Escrita:} Inércia da memória quando \texttt{mem\_write = 0}.
          \item \textbf{Leitura Protegida (\texttt{lw}):} Saída forçada a zero quando \texttt{mem\_read = 0}, evitando propagação de ruído.
          \item \textbf{Múltiplos Endereços:} Teste com gravações e leituras sequenciais.
        \end{itemize}
        \textbf{Status:} \textcolor{passgreen}{\textbf{4 / 4 cenários aprovados com 0 erros.}}
      \end{block}
    \end{column}
  \end{columns}
\end{frame}

% ------------------------------------------------------------------------------
% Slide 10: O Grande Desafio: Integração e o Algoritmo Sort
% ------------------------------------------------------------------------------
\begin{frame}{Teste Top-Level: O Desafio do Algoritmo Sort}
  \begin{block}{Cenário Real da Disciplina}
    O teste definitivo da Entrega 2 consistiu em executar o código de máquina do algoritmo de ordenação (\texttt{sort.bin}), derivado do assembly MIPS \texttt{sort.s}, comprovando a capacidade de ordenar um vetor na memória RAM.
  \end{block}

  \vspace{0.3cm}
  \textbf{Desafios Técnicos Superados Sem Alterar o Binário Fornecido:}
  \begin{itemize}\setlength{\itemsep}{6pt}
    \item \textbf{Instruções Estendidas:} O \texttt{sort.s} demandava \texttt{sll} (cálculo de deslocamento do array), \texttt{slti} (teste de índice $j < 0$) e \texttt{bne} (desvio por diferença). O hardware foi expandido com suporte nativo a essas 3 instruções.
    \item \textbf{Bypass de Endereço Base do MARS:} O montador MARS gera endereços com base em \texttt{0x00400000}. Uma máscara em nível de barramento ($\text{addr}[9:2]$) mapeia essas referências diretamente nas 256 palavras da ROM.
    \item \textbf{Bootloader de Hardware no Reset:} O código do \texttt{sort} é uma função pura (procedimento). Ao ligar o hardware, o registrador de reset inicializa automaticamente:
      $$\mathbf{\$sp} = 256 \text{ (Pilha no topo da RAM)}, \quad \mathbf{\$a0} = 0 \text{ (Ponteiro do vetor)}, \quad \mathbf{\$a1} = 5 \text{ (Tamanho do vetor)}$$
  \end{itemize}
\end{frame}

% ------------------------------------------------------------------------------
% Slide 11: Resultados da Validação Top-Level: Ordenação do Vetor
% ------------------------------------------------------------------------------
\begin{frame}{Resultado da Validação Top-Level: Antes e Depois na RAM}
  A simulação do processador completo (\texttt{tb\_mips\_monociclo\_top.v}) comprovou a execução autônoma de centenas de ciclos de clock, manipulando laços aninhados e realizando chamadas \texttt{jal swap} e retornos \texttt{jr \$ra}.

  \vspace{0.3cm}
  \begin{columns}[T]
    \begin{column}{0.48\textwidth}
      \begin{alertblock}{Estado Inicial da Memória RAM (Tempo 0)}
        \centering
        \textbf{Vetor Desordenado na RAM:} \\
        \vspace{0.2cm}
        \small
        \begin{tabular}{ccc}
          \toprule
          \textbf{Endereço} & \textbf{Índice} & \textbf{Valor Armazenado} \\
          \midrule
          \texttt{0x00} & \texttt{ram[0]} & \textbf{88} \\
          \texttt{0x04} & \texttt{ram[1]} & \textbf{15} \\
          \texttt{0x08} & \texttt{ram[2]} & \textbf{42} \\
          \texttt{0x0C} & \texttt{ram[3]} & \textbf{7} \\
          \texttt{0x10} & \texttt{ram[4]} & \textbf{99} \\
          \bottomrule
        \end{tabular}
      \end{alertblock}
    \end{column}

    \begin{column}{0.48\textwidth}
      \begin{exampleblock}{Estado Final da Memória RAM (Fim do Sort)}
        \centering
        \textbf{Vetor Perfeitamente Ordenado:} \\
        \vspace{0.2cm}
        \small
        \begin{tabular}{ccc}
          \toprule
          \textbf{Endereço} & \textbf{Índice} & \textbf{Valor Armazenado} \\
          \midrule
          \texttt{0x00} & \texttt{ram[0]} & \textcolor{passgreen}{\textbf{7}} \\
          \texttt{0x04} & \texttt{ram[1]} & \textcolor{passgreen}{\textbf{15}} \\
          \texttt{0x08} & \texttt{ram[2]} & \textcolor{passgreen}{\textbf{42}} \\
          \texttt{0x0C} & \texttt{ram[3]} & \textcolor{passgreen}{\textbf{88}} \\
          \texttt{0x10} & \texttt{ram[4]} & \textcolor{passgreen}{\textbf{99}} \\
          \bottomrule
        \end{tabular}
      \end{exampleblock}
    \end{column}
  \end{columns}

  \vspace{0.3cm}
  \centering
  \textbf{Conclusão:} O algoritmo de ordenação rodou integralmente no MIPS monociclo, executando leitura, comparação e troca (*swap*) com escrita síncrona na mesma região da memória.
\end{frame}

% ------------------------------------------------------------------------------
% Slide 12: Inspeção Temporal no GTKWave
% ------------------------------------------------------------------------------
\begin{frame}{Inspeção Temporal de Sinais no GTKWave}
  \begin{block}{Geração e Análise de Formas de Onda (Arquivos \texttt{.vcd})}
    Todos os testbenches geram dumps VCD em \texttt{build/}, visualizáveis via comandos \texttt{make waves} e \texttt{make waves-top}.
  \end{block}

  \vspace{0.3cm}
  \begin{columns}[T]
    \begin{column}{0.48\textwidth}
      \textbf{Eventos Críticos Rastreados no Top-Level:}
      \begin{itemize}\setlength{\itemsep}{5pt}
        \item \textbf{Evolução do PC:} Transição contínua $\text{PC} \leftarrow \text{PC}+4$, interceptada por saltos para rótulos de laço (\texttt{for1tst}, \texttt{for2tst}).
        \item \textbf{Chamada com Link (\texttt{jal}):} Barramento \texttt{reg\_write\_data} recebendo $\text{PC}+4$ e endereçamento de escrita em \texttt{write\_reg = 5'd31} (\$ra).
        \item \textbf{Retorno (\texttt{jr}):} Barramento \texttt{next\_pc} recebendo instantaneamente o conteúdo de \$ra através do multiplexador \texttt{mux\_jr\_inst}.
      \end{itemize}
    \end{column}

    \begin{column}{0.48\textwidth}
      \begin{exampleblock}{Evidências de Escrita na RAM}
        \begin{itemize}\setlength{\itemsep}{5pt}
          \item \textbf{Pulsos de \texttt{mem\_write}:} Ativação exclusiva nas instruções \texttt{sw} da sub-rotina \texttt{swap} e salvamento da pilha (\$sp).
          \item \textbf{Pilha Operacional:} Ajuste do ponteiro de pilha via \texttt{addi \$sp, \$sp, -20} preservando os registradores salvos (\$s0-\$s3 e \$ra).
          \item \textbf{Sinais da ULA:} Transição alternada da flag \texttt{zero} e resultados da comparação \texttt{slt}/\texttt{slti}.
        \end{itemize}
      \end{exampleblock}
    \end{column}
  \end{columns}
\end{frame}

% ------------------------------------------------------------------------------
% Slide 13: Cumprimento das Rubricas da Disciplina
% ------------------------------------------------------------------------------
\begin{frame}{Conformidade com as Rubricas de Avaliação}
  \centering
  \small
  \begin{tabular}{lp{8.2cm}c}
    \toprule
    \textbf{Rubrica} & \textbf{Requisito Oficial} & \textbf{Status} \\
    \midrule
    \textbf{Rubrica 1} & Configuração do ambiente de desenvolvimento e síntese no Quartus Prime & \textcolor{passgreen}{\textbf{100\% Concluído}} \\
    \addlinespace[3pt]
    \textbf{Rubrica 2} & Caminho de dados completo com suporte a Jumps (\texttt{j}, \texttt{jal}, \texttt{jr}) & \textcolor{passgreen}{\textbf{100\% Concluído}} \\
    \addlinespace[3pt]
    \textbf{Rubrica 3} & Validação de \textbf{todos} os componentes via testbenches individuais e top-level & \textcolor{passgreen}{\textbf{100\% Concluído}} \\
    \addlinespace[3pt]
    \textbf{Rubrica 4} & Código 100\% parametrizado (\texttt{parameter DATA\_WIDTH = 32}, etc.) & \textcolor{passgreen}{\textbf{100\% Concluído}} \\
    \addlinespace[3pt]
    \textbf{Rubrica 5} & Apresentação em vídeo (YouTube) e relatório técnico padrão IEEE & \textcolor{passgreen}{\textbf{100\% Concluído}} \\
    \bottomrule
  \end{tabular}
\end{frame}

% ------------------------------------------------------------------------------
% Slide 14: Conclusão e Encerramento
% ------------------------------------------------------------------------------
\begin{frame}{Conclusão e Encerramento}
  \begin{block}{Síntese dos Resultados Alcançados}
    \begin{itemize}\setlength{\itemsep}{6pt}
      \item \textbf{Processador MIPS Monociclo Operacional:} Implementado em Verilog HDL sintetizável e parametrizado.
      \item \textbf{Suíte de Testes Exaustiva:} 7 testbenches unitários e de integração aprovados com 100\% de sucesso.
      \item \textbf{Validação Prática com Algoritmo Real:} Execução autônoma do algoritmo de ordenação \texttt{sort.bin}, comprovando leitura, desvios e escrita síncrona na memória RAM.
      \item \textbf{Pronto para o FPGA:} Projeto compilado e sintetizado com 0 erros no Intel Quartus Prime para a placa DE10-Lite.
    \end{itemize}
  \end{block}

  \vspace{0.4cm}
  \centering
  {\Large \textbf{Obrigado pela atenção!}} \\
  \vspace{0.2cm}
  \footnotesize
  \textbf{Equipe:} José Dercy Vieira da Silva Filho $\cdot$ Vinicius Menezes \\
  Universidade Federal do Amazonas -- UFAM / FEEC
\end{frame}

\end{document}
